\documentclass[11pt,aspectratio=169]{beamer}
\usetheme{SimplePlus}

\usepackage[T1]{fontenc}
\usepackage[utf8]{inputenc}
\usepackage{amsmath,amssymb,mathtools}
\usepackage{booktabs}
\usepackage{array}
\usepackage{appendixnumberbeamer}
\usepackage{hyperref}

\newcommand{\E}{\mathbb{E}}
\newcommand{\Prob}{\mathbb{P}}
\newcommand{\Var}{\mathrm{Var}}
\newcommand{\Cov}{\mathrm{Cov}}
\newcommand{\se}{\mathrm{SE}}
\newcommand{\CI}{\mathrm{CI}}
\newcommand{\ATE}{\mathrm{ATE}}
\newcommand{\att}{\mathrm{ATT}}
\newcommand{\Normal}{\mathcal{N}}
\newcommand{\indep}{\perp\!\!\!\perp}
\newcommand{\derivbutton}[1]{\vspace{0.4em}\hfill\hyperlink{app:#1}{\beamerbutton{Appendix details}}}
\newcommand{\backbutton}[1]{\vspace{0.4em}\hfill\hyperlink{main:#1}{\beamerbutton{Back to main slide}}}

\title{Applied Econometrics: Session 3}
\subtitle{Stats Review: The Most Dangerous Equation}
\author{PhD Intensive Course}
\institute{Lecture 3}
\date{\today}

\begin{document}

\begin{frame}
  \titlepage
\end{frame}

\begin{frame}{Why This Lecture Matters}
  Lectures 1--2 focused on causal meaning: potential outcomes, experiments, selection bias, and regression as a CEF approximation.
  \vspace{0.4em}

  Today we add the missing statistical layer:
  \begin{block}{Core question}
    Once we have an estimate, how much should we trust its numerical value?
  \end{block}

  \begin{itemize}
    \item A correct design can still produce a noisy estimate.
    \item Small samples generate surprisingly extreme averages.
    \item Standard errors, confidence intervals, tests, and p-values are tools for describing that noise.
  \end{itemize}
\end{frame}

\begin{frame}{Roadmap for Today}
  \tableofcontents
\end{frame}

\section{Sampling Variation}

\begin{frame}{The Three Objects Students Must Separate}
  \begin{center}
  \begin{tabular}{lll}
    \toprule
    Object & Meaning & Example \\
    \midrule
    Estimand & Target in the population & \(\ATE=\E[Y_1-Y_0]\) \\
    Estimator & Rule applied to data & \(\bar{Y}_1-\bar{Y}_0\) \\
    Estimate & Number from this sample & \(4.2\) test points \\
    \bottomrule
  \end{tabular}
  \end{center}
  \vspace{0.6em}
  \begin{block}{Beginner mistake}
    Treating the estimate as if it were the estimand. The estimate is one noisy draw from a sampling process.
  \end{block}
\end{frame}

\begin{frame}{One Sample Is Not the Whole Story}
  Imagine a population of students with a true average score \(\mu=70\).
  \begin{itemize}
    \item Draw one classroom of 25 students: the sample average might be 66.
    \item Draw another classroom of 25 students: the sample average might be 73.
    \item Neither classroom average is ``the truth''; each is a noisy estimate.
  \end{itemize}
  \vspace{0.3em}
  \begin{block}{Sampling variation}
    If we repeated the same sampling procedure many times, the estimate would vary from sample to sample.
  \end{block}
\end{frame}

\begin{frame}{The School-Ranking Trap}
  Suppose we rank schools by average test score.
  \begin{itemize}
    \item Some schools have 20 tested students.
    \item Others have 500 tested students.
    \item Even if all schools are equally good, the small schools will have much noisier averages.
  \end{itemize}
  \vspace{0.3em}
  \begin{block}{Policy danger}
    The schools at the top of the ranking may look special partly because they had few students and a lucky draw.
  \end{block}
\end{frame}

\begin{frame}{Look at Both Tails}
  The mistake is clearer when we inspect both ends of the ranking.
  \begin{itemize}
    \item Small schools are overrepresented among the top performers.
    \item But small schools are also overrepresented among the bottom performers.
    \item This is a variance pattern, not evidence that small schools are automatically better or worse.
  \end{itemize}
  \vspace{0.3em}
  \begin{block}{Diagnostic}
    When a small group appears often among the winners, check whether it also appears often among the losers.
  \end{block}
\end{frame}

\begin{frame}{The Most Dangerous Equation}
  \hypertarget{main:dangerous-equation}{}
  For a sample mean,
  \[
    \se(\bar{Y})=\frac{\sigma}{\sqrt{n}}.
  \]
  \begin{itemize}
    \item \(\sigma\): individual-level standard deviation.
    \item \(n\): number of observations used in the average.
    \item \(\se(\bar{Y})\): standard deviation of the estimate across repeated samples.
  \end{itemize}
  \begin{block}{Translation}
    More data makes averages more stable, but precision improves with \(\sqrt{n}\), not with \(n\).
  \end{block}
  \derivbutton{mean-variance}
\end{frame}

\begin{frame}{Why Precision Improves Slowly}
  If individual scores have \(\sigma=12\):
  \begin{center}
  \begin{tabular}{rcc}
    \toprule
    Sample size \(n\) & \(\se(\bar{Y})=12/\sqrt{n}\) & Interpretation \\
    \midrule
    25 & 2.40 & very noisy average \\
    100 & 1.20 & half as noisy \\
    400 & 0.60 & half again \\
    1600 & 0.30 & half again \\
    \bottomrule
  \end{tabular}
  \end{center}
  \vspace{0.4em}
  \textbf{Rule:} to cut the standard error in half, you need four times as many observations.
\end{frame}

\begin{frame}{Standard Deviation vs Standard Error}
  \begin{center}
  \begin{tabular}{p{0.42\linewidth}p{0.42\linewidth}}
    \toprule
    Standard deviation & Standard error \\
    \midrule
    Spread of individual outcomes & Spread of the estimator \\
    Describes students, workers, firms & Describes \(\bar{Y}\), \(\hat{\beta}\), \(\widehat{\ATE}\) \\
    Usually does not vanish with \(n\) & Shrinks as sample size grows \\
    Answers: how different are units? & Answers: how precise is the estimate? \\
    \bottomrule
  \end{tabular}
  \end{center}
  \vspace{0.4em}
  \begin{block}{Practical warning}
    A dataset can have very dispersed outcomes and still estimate the mean precisely if \(n\) is large.
  \end{block}
\end{frame}

\begin{frame}{Estimators Are Random Variables}
  \hypertarget{main:estimator-rv}{}
  Before seeing data, \(\bar{Y}\) is random because the sample is random.
  \[
    \bar{Y}
    =
    \frac{1}{n}\sum_{i=1}^n Y_i.
  \]
  \begin{itemize}
    \item Different samples produce different values of \(\bar{Y}\).
    \item The sampling distribution describes those possible values.
    \item The standard error is the standard deviation of that sampling distribution.
  \end{itemize}
  \begin{block}{Mental model}
    Think of the estimate in your paper as one point from a cloud of estimates that could have happened.
  \end{block}
  \derivbutton{sampling-distribution}
\end{frame}

\begin{frame}{Bias and Variance Are Different Problems}
  \[
    \text{Estimation error}
    =
    \underbrace{\E[\hat{\theta}]-\theta}_{\text{bias}}
    +
    \underbrace{\hat{\theta}-\E[\hat{\theta}]}_{\text{sampling noise}}.
  \]
  \begin{itemize}
    \item Design assumptions target bias: selection, omitted variables, bad controls.
    \item Standard errors target sampling noise: how much \(\hat{\theta}\) moves across samples.
    \item A tiny standard error does not make a biased comparison causal.
    \item A randomized experiment can be unbiased but still imprecise.
  \end{itemize}
\end{frame}

\begin{frame}{Python Demonstration for This Lecture}
  The file
  \[
    \texttt{lectures/code/lecture-3-stats-review.py}
  \]
  simulates two examples:
  \begin{itemize}
    \item school rankings where small schools appear in both tails because their averages are noisy;
    \item randomized treatment-control studies where the true treatment effect is fixed but estimates vary across samples.
  \end{itemize}
  \vspace{0.3em}
  Run:
  \[
    \texttt{python lectures/code/lecture-3-stats-review.py}
  \]
  Optional plot export:
  \[
    \texttt{python lectures/code/lecture-3-stats-review.py --save-plots}
  \]
\end{frame}

\section{Standard Errors}

\begin{frame}{The Standard Error of a Mean}
  \hypertarget{main:se-mean}{}
  In practice \(\sigma\) is unknown, so we estimate it with the sample standard deviation:
  \[
    s
    =
    \sqrt{\frac{1}{n-1}\sum_{i=1}^n (Y_i-\bar{Y})^2}.
  \]
  Then
  \[
    \widehat{\se}(\bar{Y})
    =
    \frac{s}{\sqrt{n}}.
  \]
  \begin{itemize}
    \item Large \(s\): outcomes are noisy, so the average is harder to estimate.
    \item Large \(n\): averaging cancels out more individual noise.
  \end{itemize}
  \derivbutton{sample-se}
\end{frame}

\begin{frame}{A Small Numerical Example}
  Suppose average exam score in one class is \(\bar{Y}=72\), and \(s=12\).
  \begin{center}
  \begin{tabular}{rcc}
    \toprule
    Number of students & \(\widehat{\se}(\bar{Y})\) & Rough 95\% margin \\
    \midrule
    25 & \(12/\sqrt{25}=2.40\) & \(1.96\times 2.40=4.70\) \\
    100 & \(12/\sqrt{100}=1.20\) & \(1.96\times 1.20=2.35\) \\
    400 & \(12/\sqrt{400}=0.60\) & \(1.96\times 0.60=1.18\) \\
    \bottomrule
  \end{tabular}
  \end{center}
  \vspace{0.4em}
  Same mean, same individual dispersion, very different precision.
\end{frame}

\begin{frame}{Treatment-Control Difference in Means}
  \hypertarget{main:ate-estimator}{}
  In a randomized experiment with binary treatment \(D_i\):
  \[
    \widehat{\ATE}
    =
    \bar{Y}_1-\bar{Y}_0,
  \]
  where
  \[
    \bar{Y}_1=\frac{1}{n_1}\sum_{D_i=1}Y_i,
    \qquad
    \bar{Y}_0=\frac{1}{n_0}\sum_{D_i=0}Y_i.
  \]
  \begin{itemize}
    \item Randomization gives the causal interpretation.
    \item The standard error tells us how noisy this causal estimate is.
  \end{itemize}
  \derivbutton{two-sample-se}
\end{frame}

\begin{frame}{Standard Error of a Difference}
  \hypertarget{main:ate-se}{}
  If treatment and control samples are independent,
  \[
    \widehat{\se}(\bar{Y}_1-\bar{Y}_0)
    =
    \sqrt{\frac{s_1^2}{n_1}+\frac{s_0^2}{n_0}}.
  \]
  \begin{itemize}
    \item Each group mean is noisy.
    \item The variance of the difference adds the two group variances.
    \item The noisier or smaller group can dominate the uncertainty.
  \end{itemize}
  \begin{block}{Causal translation}
    Even with random assignment, the estimated effect can look too large, too small, or even have the wrong sign in a small sample.
  \end{block}
  \derivbutton{two-sample-se}
\end{frame}

\begin{frame}{Unequal Group Sizes Waste Precision}
  Suppose \(s_1=s_0=12\) and total sample size is \(200\).
  \begin{center}
  \begin{tabular}{rcc}
    \toprule
    \(n_1\) treated & \(n_0\) control & \(\widehat{\se}(\bar{Y}_1-\bar{Y}_0)\) \\
    \midrule
    100 & 100 & 1.70 \\
    50 & 150 & 1.96 \\
    20 & 180 & 2.83 \\
    \bottomrule
  \end{tabular}
  \end{center}
  \vspace{0.3em}
  \begin{block}{Design lesson}
    For a fixed total sample and equal variances, balanced treatment-control groups are most precise.
  \end{block}
\end{frame}

\begin{frame}{What Standard Errors Do Not Do}
  Standard errors are essential, but they are not magic.
  \begin{itemize}
    \item They do not remove selection bias.
    \item They do not validate a bad control strategy.
    \item They do not prove the outcome was measured correctly.
    \item They do not solve spillovers, attrition, or noncompliance.
  \end{itemize}
  \begin{block}{Good applied work needs both}
    Identification answers: is the estimate centered on the causal target? Inference answers: how uncertain is the estimate around its center?
  \end{block}
\end{frame}

\section{Confidence Intervals}

\begin{frame}{The CLT as the Workhorse}
  \hypertarget{main:clt}{}
  Under common regularity conditions,
  \[
    \frac{\bar{Y}-\mu}{\se(\bar{Y})}
    \approx
    \Normal(0,1).
  \]
  \begin{itemize}
    \item The sample mean is approximately normal in large samples.
    \item The approximation often works better for averages than for individual data.
    \item We use this to translate standard errors into intervals and tests.
  \end{itemize}
  \derivbutton{clt-ci}
\end{frame}

\begin{frame}{A 95\% Confidence Interval}
  \hypertarget{main:ci-formula}{}
  For an estimate \(\hat{\theta}\),
  \[
    \hat{\theta}
    \pm
    1.96\,\widehat{\se}(\hat{\theta})
  \]
  is the usual large-sample 95\% confidence interval.
  \begin{itemize}
    \item Center: the estimate from this sample.
    \item Radius: about two standard errors.
    \item Wider interval: less precise estimate.
    \item Narrower interval: more precise estimate.
  \end{itemize}
  \derivbutton{clt-ci}
\end{frame}

\begin{frame}{Treatment Effect Interval: Example}
  Suppose a randomized tutoring program produces:
  \[
    \bar{Y}_1-\bar{Y}_0=4.0,
    \qquad
    \widehat{\se}=1.5.
  \]
  Then
  \[
    95\%~\CI
    =
    4.0\pm 1.96(1.5)
    =
    [1.06,\;6.94].
  \]
  \begin{block}{Substantive reading}
    The data are compatible with a modest positive effect and a large positive effect, but not with exactly zero under the usual 95\% rule.
  \end{block}
\end{frame}

\begin{frame}{What Does 95\% Confidence Mean?}
  \hypertarget{main:ci-interpretation}{}
  Frequentist confidence is about the procedure, not about a fixed interval.
  \begin{itemize}
    \item Imagine repeating the study many times.
    \item Each time, compute \(\hat{\theta}\pm1.96\,\widehat{\se}\).
    \item About 95\% of those intervals would cover the true parameter.
  \end{itemize}
  \begin{block}{Careful language}
    After seeing one interval, the true parameter is fixed. The interval either covers it or it does not.
  \end{block}
  \derivbutton{clt-ci}
\end{frame}

\begin{frame}{Common Confidence-Interval Mistakes}
  \begin{itemize}
    \item Wrong: ``There is a 95\% probability that this exact interval contains the true parameter.''
    \item Better: ``This method covers the true parameter in about 95\% of repeated samples.''
    \item Wrong: ``All values inside the interval are equally plausible.''
    \item Better: ``Values near the estimate are usually more compatible with the data than values near the edge.''
    \item Wrong: ``A wide interval means the design is not causal.''
    \item Better: ``A wide interval means the estimate is imprecise.''
  \end{itemize}
\end{frame}

\begin{frame}{Do Not Compare Two Separate CIs}
  \hypertarget{main:overlap}{}
  To test whether treatment and control means differ, use the CI for the difference:
  \[
    (\bar{Y}_1-\bar{Y}_0)
    \pm
    1.96
    \sqrt{\frac{s_1^2}{n_1}+\frac{s_0^2}{n_0}}.
  \]
  \begin{itemize}
    \item Separate CIs for \(\bar{Y}_1\) and \(\bar{Y}_0\) answer separate questions.
    \item Overlap of those separate CIs is not the right test for the treatment effect.
    \item The correct uncertainty is attached to the estimator you report.
  \end{itemize}
  \derivbutton{overlap-diff}
\end{frame}

\begin{frame}{Interval Width Is Information}
  Compare two estimates of the same causal effect:
  \begin{center}
  \begin{tabular}{lcc}
    \toprule
    Study & Estimate & 95\% CI \\
    \midrule
    Small pilot & 5.0 & \([-3.0,\;13.0]\) \\
    Large RCT & 5.0 & \([3.5,\;6.5]\) \\
    \bottomrule
  \end{tabular}
  \end{center}
  \vspace{0.4em}
  Same point estimate, very different evidence.
  \begin{itemize}
    \item The pilot is compatible with harm, no effect, or large benefit.
    \item The large RCT pins down a positive effect much more tightly.
  \end{itemize}
\end{frame}

\section{Hypothesis Tests and P-values}

\begin{frame}{Hypothesis Testing: The Logic}
  \hypertarget{main:hyp-test}{}
  A hypothesis test begins by assuming a benchmark is true.
  \[
    H_0:\theta=\theta_0.
  \]
  Then ask:
  \begin{block}{Testing question}
    If \(H_0\) were true, would an estimate this far from \(\theta_0\) be surprising?
  \end{block}
  \begin{itemize}
    \item If very surprising, reject \(H_0\).
    \item If not very surprising, do not reject \(H_0\).
    \item ``Do not reject'' is not the same as ``prove true.''
  \end{itemize}
  \derivbutton{z-test}
\end{frame}

\begin{frame}{The Test Statistic}
  \hypertarget{main:z-stat}{}
  For a large-sample normal test,
  \[
    z
    =
    \frac{\hat{\theta}-\theta_0}
         {\widehat{\se}(\hat{\theta})}.
  \]
  \begin{itemize}
    \item Numerator: distance from the null value.
    \item Denominator: typical sampling noise.
    \item A large absolute \(z\) means the estimate is far from the null relative to its noise.
  \end{itemize}
  \begin{block}{95\% two-sided rule}
    Reject \(H_0\) at the 5\% level if \(|z|>1.96\).
  \end{block}
  \derivbutton{z-test}
\end{frame}

\begin{frame}{P-values}
  \hypertarget{main:pvalue}{}
  The p-value is:
  \[
    \Prob(\text{test statistic at least this extreme}\mid H_0).
  \]
  For a two-sided normal test:
  \[
    p
    =
    2\left[1-\Phi\left(|z|\right)\right].
  \]
  \begin{itemize}
    \item Small p-value: this estimate would be unusual if \(H_0\) were true.
    \item Large p-value: this estimate is not unusual under \(H_0\).
  \end{itemize}
  \derivbutton{pvalue-calc}
\end{frame}

\begin{frame}{The Most Common P-value Mistake}
  A p-value is not
  \[
    \Prob(H_0\mid \text{data}).
  \]
  It is closer to
  \[
    \Prob(\text{data as or more extreme}\mid H_0).
  \]
  \begin{block}{Plain English}
    A p-value asks how surprising the estimate is under the null. It does not tell us the probability that the null is true.
  \end{block}
\end{frame}

\begin{frame}{One-sided and Two-sided Tests}
  \begin{itemize}
    \item Two-sided alternative:
    \[
      H_1:\theta\neq\theta_0.
    \]
    Reject for estimates far above or far below the null.
    \item One-sided alternative:
    \[
      H_1:\theta>\theta_0
      \quad\text{or}\quad
      H_1:\theta<\theta_0.
    \]
    Reject only in the pre-specified direction.
  \end{itemize}
  \begin{block}{Applied convention}
    Use two-sided tests unless the one-sided direction was justified before seeing the data.
  \end{block}
\end{frame}

\begin{frame}{Statistical Significance Is Not Economic Importance}
  \begin{center}
  \begin{tabular}{lcc}
    \toprule
    Result & Estimate & Standard error \\
    \midrule
    Huge dataset & 0.03 & 0.01 \\
    Small study & 4.00 & 2.50 \\
    \bottomrule
  \end{tabular}
  \end{center}
  \vspace{0.4em}
  \begin{itemize}
    \item The first effect may be statistically significant but economically tiny.
    \item The second effect may be substantively large but statistically imprecise.
    \item Applied interpretation must discuss magnitude, uncertainty, and design.
  \end{itemize}
\end{frame}

\begin{frame}{Failure to Reject Is Not Evidence of Zero}
  Suppose \(\hat{\theta}=4.0\) and \(\widehat{\se}=3.0\).
  \[
    z=\frac{4.0}{3.0}=1.33,
    \qquad
    95\%~\CI=[-1.88,\;9.88].
  \]
  \begin{itemize}
    \item We cannot reject zero at the 5\% level.
    \item But the interval also contains effects large enough to matter.
    \item The study may simply be underpowered.
  \end{itemize}
  \begin{block}{Better sentence}
    ``The estimate is too imprecise to distinguish zero from meaningful positive effects.''
  \end{block}
\end{frame}

\begin{frame}{Power: Why Small Studies Miss Effects}
  Power is the probability of rejecting \(H_0\) when a particular alternative is true.
  \begin{itemize}
    \item Small \(n\) gives large standard errors.
    \item Large standard errors make \(z=(\hat{\theta}-\theta_0)/\widehat{\se}\) smaller in absolute value.
    \item Low-power studies often produce unstable estimates and inconsistent conclusions.
  \end{itemize}
  \begin{block}{Applied implication}
    A noisy null result is not strong evidence that the causal effect is zero.
  \end{block}
\end{frame}

\begin{frame}{Multiple Testing and Specification Search}
  If you test many outcomes, subgroups, or specifications:
  \begin{itemize}
    \item some small p-values can arise by chance;
    \item the most impressive estimate is often selected after looking at many estimates;
    \item standard errors for one chosen regression do not account for all the searching.
  \end{itemize}
  \begin{block}{Practical discipline}
    Pre-specify main outcomes when possible, report the full family of results, and be honest about exploratory analyses.
  \end{block}
\end{frame}

\section{Uncertainty and Causal Estimates}

\begin{frame}{Design Credibility vs Statistical Precision}
  \begin{center}
  \begin{tabular}{p{0.38\linewidth}p{0.45\linewidth}}
    \toprule
    Problem & Main question \\
    \midrule
    Identification & Is the estimator centered on the causal estimand? \\
    Inference & How dispersed is the estimator around that center? \\
    Measurement & Is the outcome/treatment recorded correctly? \\
    External validity & Does this population answer the policy question? \\
    \bottomrule
  \end{tabular}
  \end{center}
  \vspace{0.4em}
  \begin{block}{Course habit}
    Never use the word ``causal'' just because a p-value is small.
  \end{block}
\end{frame}

\begin{frame}{Randomization Removes Selection, Not Luck}
  \hypertarget{main:randomization-noise}{}
  In an RCT:
  \[
    D_i\indep(Y_{1i},Y_{0i})
  \]
  gives unbiased comparison of treatment and control groups.
  \vspace{0.3em}

  But in a finite sample:
  \[
    \bar{Y}_1-\bar{Y}_0
  \]
  still varies across random assignments and random samples.
  \begin{block}{Lesson}
    Randomization solves the identification problem in expectation. Standard errors describe the remaining finite-sample uncertainty.
  \end{block}
  \derivbutton{randomization-variance}
\end{frame}

\begin{frame}{Observational Causal Designs Need More Care}
  In observational studies, uncertainty often has additional structure:
  \begin{itemize}
    \item heteroskedasticity: outcome variance differs across units;
    \item clustering: units in the same school, firm, village, or year share shocks;
    \item serial correlation: repeated observations over time are not independent;
    \item generated regressors: first-stage estimates add uncertainty.
  \end{itemize}
  \begin{block}{Practical rule}
    The standard error formula must match the design and sampling process.
  \end{block}
\end{frame}

\begin{frame}{What to Report in Applied Work}
  For every main causal estimate, report:
  \begin{itemize}
    \item point estimate and units;
    \item standard error;
    \item confidence interval;
    \item sample size and group sizes;
    \item p-value if testing is relevant;
    \item the identifying assumption that gives the estimate causal meaning.
  \end{itemize}
  \begin{block}{Minimum standard}
    An effect estimate without uncertainty is not yet an empirical result.
  \end{block}
\end{frame}

\begin{frame}{Reading a Results Table}
  When you see a table, read it in this order:
  \begin{enumerate}
    \item What is the estimand? ATE, ATT, regression slope, reduced form?
    \item What design or assumption identifies it?
    \item What is the sign and magnitude?
    \item What is the standard error or confidence interval?
    \item Is the estimate precise enough to answer the economic question?
  \end{enumerate}
  \vspace{0.3em}
  Stars are last, not first.
\end{frame}

\begin{frame}{A Useful Writing Template}
  \begin{block}{Template}
    We estimate that \(X\) increases/decreases \(Y\) by \(\hat{\theta}\) units
    \((\widehat{\se}=s)\), with a 95\% confidence interval of
    \([L,U]\). The identifying variation comes from \(\ldots\).
  \end{block}
  \vspace{0.4em}
  Then interpret:
  \begin{itemize}
    \item Is \(\hat{\theta}\) economically large?
    \item Does \([L,U]\) rule out policy-relevant alternatives?
    \item What assumption would move the estimate away from the causal target?
  \end{itemize}
\end{frame}

\begin{frame}{Lecture Takeaways}
  \begin{enumerate}
    \item Small samples naturally create extreme estimates.
    \item The standard error of a mean is \(\sigma/\sqrt{n}\).
    \item Confidence intervals describe uncertainty around an estimate.
    \item Hypothesis tests ask whether data are surprising under a null.
    \item P-values are \(P(\text{data as extreme}\mid H_0)\), not \(P(H_0\mid\text{data})\).
    \item Causal interpretation requires identification; statistical precision is a separate requirement.
  \end{enumerate}
\end{frame}

\begin{frame}{Sources and Further Reading}
  \begin{itemize}
    \item Matheus Facure Alves, \emph{Causal Inference for the Brave and True}, Chapter 03:
    \href{https://matheusfacure.github.io/python-causality-handbook/03-Stats-Review-The-Most-Dangerous-Equation.html}{Stats Review: The Most Dangerous Equation}.
    \item Howard Wainer (2007), ``The Most Dangerous Equation,'' \emph{American Scientist}, 95(3), 249--256.
    \item Angrist and Pischke, \emph{Mostly Harmless Econometrics}, review of experiments and standard errors.
    \item Angrist and Pischke, \emph{Mastering 'Metrics}, discussion of experiments and regression evidence.
  \end{itemize}
  \vspace{0.3em}
  These slides paraphrase and extend the chapter for this course; the examples use self-contained simulations.
\end{frame}

\appendix

\begin{frame}{Appendix Map}
  \begin{itemize}
    \item \hyperlink{app:mean-variance}{\beamerbutton{A1: Variance of the sample mean}}
    \item \hyperlink{app:sampling-distribution}{\beamerbutton{A2: Sampling distributions}}
    \item \hyperlink{app:sample-se}{\beamerbutton{A3: Sample standard deviation and plug-in SE}}
    \item \hyperlink{app:clt-ci}{\beamerbutton{A4: CLT and confidence intervals}}
    \item \hyperlink{app:two-sample-se}{\beamerbutton{A5: Two-sample standard error}}
    \item \hyperlink{app:z-test}{\beamerbutton{A6: Hypothesis tests}}
    \item \hyperlink{app:pvalue-calc}{\beamerbutton{A7: P-value calculation}}
    \item \hyperlink{app:overlap-diff}{\beamerbutton{A8: Why overlapping CIs are not the test}}
    \item \hyperlink{app:randomization-variance}{\beamerbutton{A9: Randomization variance}}
  \end{itemize}
\end{frame}

\begin{frame}{Appendix A1: Variance of the Sample Mean}
  \hypertarget{app:mean-variance}{}
  Let \(Y_1,\ldots,Y_n\) be independent draws with
  \[
    \E[Y_i]=\mu,
    \qquad
    \Var(Y_i)=\sigma^2.
  \]
  The sample mean is
  \[
    \bar{Y}=\frac{1}{n}\sum_{i=1}^n Y_i.
  \]
  Then
  \[
    \Var(\bar{Y})
    =
    \Var\left(\frac{1}{n}\sum_{i=1}^n Y_i\right)
    =
    \frac{1}{n^2}\sum_{i=1}^n \Var(Y_i)
    =
    \frac{\sigma^2}{n}.
  \]
  Therefore
  \[
    \se(\bar{Y})=\sqrt{\Var(\bar{Y})}=\frac{\sigma}{\sqrt{n}}.
  \]
  \backbutton{dangerous-equation}
\end{frame}

\begin{frame}{Appendix A2: Sampling Distributions}
  \hypertarget{app:sampling-distribution}{}
  The population distribution describes individual \(Y_i\).
  The sampling distribution describes an estimator such as \(\bar{Y}\).
  \vspace{0.3em}

  If
  \[
    Y_i \sim \Normal(\mu,\sigma^2),
  \]
  then exactly
  \[
    \bar{Y}
    \sim
    \Normal\left(\mu,\frac{\sigma^2}{n}\right).
  \]
  For non-normal \(Y_i\), the central limit theorem gives the approximation
  \[
    \bar{Y}
    \approx
    \Normal\left(\mu,\frac{\sigma^2}{n}\right)
  \]
  in large samples.
  \backbutton{estimator-rv}
\end{frame}

\begin{frame}{Appendix A3: Sample Standard Deviation and Plug-in SE}
  \hypertarget{app:sample-se}{}
  The unknown population variance is
  \[
    \sigma^2=\E[(Y_i-\mu)^2].
  \]
  In data, replace \(\mu\) with \(\bar{Y}\):
  \[
    s^2
    =
    \frac{1}{n-1}\sum_{i=1}^n (Y_i-\bar{Y})^2.
  \]
  The denominator is \(n-1\), not \(n\), because estimating \(\bar{Y}\) consumes one degree of freedom.
  \vspace{0.3em}

  Plug this into de Moivre's formula:
  \[
    \widehat{\se}(\bar{Y})
    =
    \sqrt{\frac{s^2}{n}}
    =
    \frac{s}{\sqrt{n}}.
  \]
  \backbutton{se-mean}
\end{frame}

\begin{frame}{Appendix A4: CLT and Confidence Intervals}
  \hypertarget{app:clt-ci}{}
  The CLT implies
  \[
    Z
    =
    \frac{\bar{Y}-\mu}{\se(\bar{Y})}
    \approx
    \Normal(0,1).
  \]
  For a standard normal variable,
  \[
    \Prob(-1.96\leq Z\leq 1.96)\approx 0.95.
  \]
  Substitute \(Z\):
  \[
    \Prob\left(
      -1.96
      \leq
      \frac{\bar{Y}-\mu}{\se(\bar{Y})}
      \leq
      1.96
    \right)
    \approx 0.95.
  \]
  Rearranging gives
  \[
    \Prob\left(
      \bar{Y}-1.96\,\se(\bar{Y})
      \leq
      \mu
      \leq
      \bar{Y}+1.96\,\se(\bar{Y})
    \right)
    \approx 0.95.
  \]
  \backbutton{ci-formula}
\end{frame}

\begin{frame}{Appendix A5: Two-Sample Standard Error}
  \hypertarget{app:two-sample-se}{}
  Let treatment and control sample means be independent:
  \[
    \bar{Y}_1=\frac{1}{n_1}\sum_{D_i=1}Y_i,
    \qquad
    \bar{Y}_0=\frac{1}{n_0}\sum_{D_i=0}Y_i.
  \]
  The estimator is
  \[
    \hat{\tau}=\bar{Y}_1-\bar{Y}_0.
  \]
  Using \(\Var(A-B)=\Var(A)+\Var(B)-2\Cov(A,B)\):
  \[
    \Var(\hat{\tau})
    =
    \Var(\bar{Y}_1)+\Var(\bar{Y}_0)
    =
    \frac{\sigma_1^2}{n_1}+\frac{\sigma_0^2}{n_0},
  \]
  where independence makes the covariance term zero. Estimate with
  \[
    \widehat{\se}(\hat{\tau})
    =
    \sqrt{\frac{s_1^2}{n_1}+\frac{s_0^2}{n_0}}.
  \]
  \backbutton{ate-se}
\end{frame}

\begin{frame}{Appendix A6: Hypothesis Tests}
  \hypertarget{app:z-test}{}
  To test
  \[
    H_0:\theta=\theta_0,
  \]
  start from the approximate distribution
  \[
    \hat{\theta}
    \approx
    \Normal\left(\theta,\se(\hat{\theta})^2\right).
  \]
  Under the null,
  \[
    \frac{\hat{\theta}-\theta_0}{\se(\hat{\theta})}
    \approx
    \Normal(0,1).
  \]
  Replacing the unknown standard error with its estimate gives
  \[
    z
    =
    \frac{\hat{\theta}-\theta_0}{\widehat{\se}(\hat{\theta})}.
  \]
  For a two-sided 5\% test, reject when \(|z|>1.96\).
  \backbutton{z-stat}
\end{frame}

\begin{frame}{Appendix A7: P-value Calculation}
  \hypertarget{app:pvalue-calc}{}
  Let \(Z\sim\Normal(0,1)\), and suppose the observed test statistic is \(z_{\text{obs}}\).
  \vspace{0.3em}

  For a two-sided test,
  \[
    p
    =
    \Prob(|Z|\geq |z_{\text{obs}}|)
    =
    2\Prob(Z\geq |z_{\text{obs}}|)
    =
    2\left[1-\Phi(|z_{\text{obs}}|)\right].
  \]
  Example:
  \[
    z_{\text{obs}}=2.3
    \quad\Rightarrow\quad
    p=2[1-\Phi(2.3)]\approx 0.021.
  \]
  This is small because values beyond \(\pm 2.3\) are rare under a standard normal null distribution.
  \backbutton{pvalue}
\end{frame}

\begin{frame}{Appendix A8: Why Overlapping CIs Are Not the Test}
  \hypertarget{app:overlap-diff}{}
  Suppose two independent means have estimates \(\bar{Y}_1,\bar{Y}_0\) and standard errors \(\se_1,\se_0\).
  Separate 95\% intervals are
  \[
    \bar{Y}_1\pm 1.96\se_1,
    \qquad
    \bar{Y}_0\pm 1.96\se_0.
  \]
  But the treatment effect is
  \[
    \hat{\tau}=\bar{Y}_1-\bar{Y}_0,
  \]
  whose standard error is
  \[
    \se(\hat{\tau})
    =
    \sqrt{\se_1^2+\se_0^2}.
  \]
  The correct 95\% interval is
  \[
    \hat{\tau}
    \pm
    1.96\sqrt{\se_1^2+\se_0^2}.
  \]
  This object can exclude zero even when the two separate mean intervals overlap.
  \backbutton{overlap}
\end{frame}

\begin{frame}{Appendix A9: Randomization Variance}
  \hypertarget{app:randomization-variance}{}
  With potential outcomes
  \[
    Y_i=D_iY_{1i}+(1-D_i)Y_{0i},
  \]
  the difference in means is
  \[
    \hat{\tau}
    =
    \bar{Y}_{1,\text{treated}}
    -
    \bar{Y}_{0,\text{control}}.
  \]
  Under random assignment, treatment status is independent of potential outcomes, so
  \[
    \E[\hat{\tau}]=\E[Y_{1i}-Y_{0i}]
  \]
  for the relevant experimental population.
  \vspace{0.3em}

  But the realized treated and control groups are finite random partitions, so
  \[
    \Var(\hat{\tau})
    \approx
    \frac{\sigma_1^2}{n_1}
    +
    \frac{\sigma_0^2}{n_0}.
  \]
  Randomization centers the estimator correctly; it does not make the variance zero.
  \backbutton{randomization-noise}
\end{frame}

\end{document}
